\documentclass{article}
\usepackage{
  ../sty/hott,
  ../sty/mathwidth,
  ../sty/infer}
% \usepackage{quiver}
\usepackage{adjustbox}


\title{Introduction to  Homotopy Type Theory\\ \large The Judgment Index}
\author{CLC}
\date{Fall 2025}

\begin{document}

\maketitle

\section{Dependent Type Theory}

\subsection{Judgments}


\begin{definition}[Judgments.]
  Judgements in MLTT are of the form:
  \begin{enumerate}
  \item $\WfType \Gamma A$
  \item $\WfTerm \Gamma a A$
  \item $\EqType \Gamma A B$
  \item $\EqTerm \Gamma a b A$
  \end{enumerate}
\end{definition}
\begin{itemize}
\item $\mathcal{J}$ is the judgment thesis,
  ${\mathcal{J} \in \{a\co A, A ~ type, M \doteq N : \sigma, A \doteq B ~ type\}}$
\item $\Gamma \vdash \mathcal{J}$ is the judgment.
\end{itemize}

\begin{definition}[Context.]
  A \emph{context} is a finite list of variable declarations
  \[ x_1 : A_1, x_2 : A_2(x_1), x_3 : A_3(x_1, x_2), ..., x_n : A_n (x_1, ..., x_n) \]
  satisfying the condition that for each $1 \leq k \leq n$
  we can derive the judgment
  \[ x_1 : A_1, x_2 : A_2(x_1), x_3 : A_3(x_1, x_2), ... \vdash x_n : A_n (x_1, ..., x_n) \]
\end{definition}

\subsection{Type Families}

\begin{definition}[Family of Types.]
  $B(x)$ is a family of types over $A$ in the context $\Gamma$
  \[ \WfType {\Gamma, x \co A} {B(x)} \]
\end{definition}

\begin{definition}[Section.]
  The element of type $B(x)$ in the context $\Gamma, x : A$ is called a \emph{section} of the family $B$ over $A$ in the context $\Gamma$
  \[ \WfTerm {\Gamma, x \co A} {b(x)} {B(x)} \]
  Alternatively, $b(x)$ is an element of type $B(x)$ \emph{indexed} by
  $x\co A$ in the context $\Gamma$
\end{definition}

\subsection{Inference Rules}
There are six sets of inference rules:
\begin{enumerate}
\item The formation of contexts, types and their elements~\cref{fig:formation-rules}.
\item Postulating of judgmental equality being an equivalence relation~\cref{fig:jeq-is-equivalence}.
\item Variable conversion
\item Substitution
\item Weakening
\item The generic element
\end{enumerate}

\begin{figure}[H]
\centering
\begin{gather*}
\ib{\irule[]
  {\WfType {\Gamma, x : A} {B(x)}};
  {\WfType \Gamma A}
}
\rsp
\ib{\irule[]
  {\EqType \Gamma A B};
  {\WfType \Gamma A}
}
\rsp
\ib{\irule[]
  {\EqType \Gamma A B};
  {\WfType \Gamma B}
}
\\
\ib{\irule[]
  {\WfTerm \Gamma a A};
  {\WfType \Gamma A}
}
\rsp
\ib{\irule[]
  {\EqTerm \Gamma a b A};
  {\WfTerm \Gamma a A}
}
\rsp
\ib{\irule[]
  {\EqTerm \Gamma a b A};
  {\WfTerm \Gamma b A}
}
\end{gather*}
\caption{Formation of Contexts, types and their elements}
\label{fig:formation-rules}
\end{figure}

\begin{figure}[H]
\begin{gather*}
\ib{\irule[\trule{Ty-Eq-R}]
  {\WfType \Gamma A};
  {\EqType \Gamma A A}
}
\rsp
\ib{\irule[\trule{Ty-Eq-S}]
  {\EqType \Gamma A B};
  {\EqType \Gamma B A}
}
\rsp
\ib{\irule[\trule{Ty-Eq-T}]
  {\EqType \Gamma A B}
  {\EqType \Gamma B C};
  {\EqType \Gamma A C}
}
\\
\ib{\irule[\trule{Tm-Eq-R}]
  {\WfTerm \Gamma M A};
  {\EqTerm \Gamma M M A}
}
\rsp
\ib{\irule[\trule{Tm-Eq-S}]
  {\EqTerm \Gamma M N A};
  {\EqTerm \Gamma N M A}
}
\rsp
\ib{\irule[\trule{Tm-Eq-T}]
  {\EqTerm \Gamma M N A}
  {\EqTerm \Gamma N O A};
  {\EqTerm \Gamma M O A}
}
\end{gather*}
\caption{Judgmental Equality is an equivalence relation}
\label{fig:jeq-is-equivalence}
\end{figure}

\begin{figure}[H]
\centering
\begin{gather*}
\ib{\irule[]
   {\EqType \Gamma A {A'}}
   { {\Gamma, x : A, \Delta} \vdash {\mathcal{J}}};
   { {\Gamma, x : A', \Delta} \vdash {\mathcal{J}}}
  }
\end{gather*}
\caption{Variable Conversion Rule}
\end{figure}

\begin{figure}[H]
\centering
\begin{gather*}
\ib{\irule[\trule{Weak}]
  {\Gamma, \Delta \vdash \mathcal{J}}
   {\WfType \Gamma \rho};
   {\Gamma, x\co\rho, \Delta \vdash \mathcal{J}}
 }
 \rsp
\ib{\irule[\trule{S}]
   {\Gamma,x\co A, \Delta \vdash \mathcal{J}}
   {\WfTerm \Gamma a A};
   {\Gamma,\Subst{x}{a}\Delta \vdash \Subst{x}{a}\mathcal{J}}
 }
\end{gather*}
\caption{Weakening and Substitution Rules}
\end{figure}

\begin{figure}[H]
  \centering
  \begin{gather*}
    \ib{\irule[\trule{$\delta$}]
      {\WfType \Gamma A};
      {\WfTerm {\Gamma,x\co A} x A}
    }
  \end{gather*}
  \caption{The generic elements}
  \label{fig:generic-elements}
\end{figure}

\section{Dependent Function Types}
\begin{figure}[H]
\centering
\begin{gather*}
\ib{\irule[\trule{$\Pi$}]
  {\WfType \Gamma A}
  {\WfType {\Gamma,x\co A} {B(x)}};
  {\WfType \Gamma {\PiTy {x\co A} {B(x)}}}
}
\rsp
\ib{\irule[\trule{\EqR{\Pi}}]
  {\EqType \Gamma {A} {A'}}
  {\EqType {\Gamma,x\co\sigma} {B(x)} {B'(x)}};
  % {\sigma \in \{A, A'\}};
  {\EqType \Gamma {\PiTy {x\co A} {B(x)}} {\PiTy {x\co A'} {B'(x)}}}
}
\\
\ib{\irule[\trule{\IR{\lambda\Pi}}]
  {\WfTerm {\Gamma,x\co A} {b(x)} \tau};
  {\WfTerm \Gamma {\Lam {x\co A} {b(x)}} {\PiTy {x\co A} \tau}}
}
\\
\ib{\irule[\trule{\EqR{\lambda\Pi}}]
  {\EqTerm {\Gamma,x\co A} {b(x)}{b'(x)} {B(x)}}
  {\EqType \Gamma { A}{ A'}};
  {\EqTerm \Gamma {\Lam {x\co A} {b(x)}} {\Lam {x\co A'} {b'(x)}} {\PiTy {x\co A} {B(x)}}}
}
\\
\ib{\irule[\trule{$ev$}]
  {\WfTerm \Gamma f {\PiTy {x\co A}  B(x)}};
  {\WfTerm {\Gamma, x \co A} {f (x)} {B(x)}}
}
\rsp
\ib{\irule[\trule{\EqR{ev}}]
  {\EqTerm \Gamma {f} {f'}{\PiTy {x\co\sigma}  {B(x)}}};
  {\EqTerm {\Gamma, x \co A} {f(x)} {f'(x)} {B(x)}}
}
\\
\ib{\irule[\trule{$\beta\Pi$}]
  {\WfTerm {\Gamma, x \co A} {b(x)} {B(x)}};
  {\EqTerm {\Gamma, x \co A} {(\Lam {y} {b(y)}) x} {b(x)} {B(x)}}
}
\rsp
\ib{\irule[\trule{$\eta\Pi$}]
  {\WfTerm \Gamma {f} {\PiTy {x\co\sigma}  {B(x)}}};
  {\EqTerm {\Gamma} {\Lam x {f(x)}} {f} {\PiTy {x\co\sigma}  {B(x)}}}
}
\end{gather*}
\caption{Dependent Function Space}
\end{figure}

\subsection{Ordinary Function types}
They are special cases of dependent function types where the binder
does not appear free in the term.
\[ A \to B :=  \PiTy {x\co A} B \]

\begin{definition}[Identity Function]
 \[\texttt{id}_A : A \to A \]
\end{definition}

\section{Dependent Sum}

\begin{figure}[H]
\centering
\begin{gather*}
\ib{\irule[\trule{\FR\Sigma}]
  {\WfType \Gamma \sigma}
  {\WfType {\Gamma,x\co\sigma} \tau};
  {\WfType \Gamma {\SigmaTy {x\co\sigma} \tau}}
}
\rsp
\ib{\irule[\trule{\EqR\Sigma}]
  {\EqType \Gamma {\sigma_1} {\sigma_2}}
  {\EqType {\Gamma,x\co\sigma} {\tau_1} {\tau_2}}
  {\sigma \in \{\sigma_1, \sigma_2\}};
  {\EqType \Gamma {\Sigma x\co\sigma_1.\tau_1} {\SigmaTy{x\co\sigma_2} {\tau_2}}}
}
\\
\ib{\irule[\trule{\IR\Sigma}]
  {\WfTerm {\Gamma} M \sigma}
  {\WfTerm \Gamma N {\Subst{x}{M}\tau}};
  {\WfTerm \Gamma {\SigmaPair {\tau[x\co\sigma]} M N} {\SigmaTy {x\co\sigma} \tau}}
}
\\
\ib{\irule[\trule{\EqR{\eta\Sigma}}]
  {\EqTerm {\Gamma} {M_1}{M_2} \sigma_1}
  {\EqTerm {\Gamma,x\co\sigma_1} {N_1}{N_2} \tau_1}
  {\EqType \Gamma {\sigma_1}{\sigma_2}}
  {\EqType {\Gamma,x\co\sigma_1} {\tau_1} {\tau_2}};
  {\EqTerm \Gamma {\SigmaPair {\tau_1[x\co\sigma_1]} {M_1} {N_1}} {\SigmaPair {\tau_2[x\co\sigma_2]} {M_2} {N_2}} {\SigmaTy {x\co\sigma_1} \tau_1}}
}
\\
\ib{\irule[\trule{\ER\Sigma}]
  {\WfType {\Gamma, z\co\SigmaTy {x\co\sigma} \tau} {\rho}}
  {\WfTerm {\Gamma, x\co\sigma, y\co\tau} H {\Subst{z}{\SigmaPair {\tau[x\co\sigma]} x y}\rho}}
  {\WfTerm \Gamma M {\SigmaTy {x\co\sigma} \tau}};
  {\EqTerm {\Gamma}{\RecSigma {\rho[z:\SigmaTy{x\co\sigma}{\tau}]} {H[x\co\sigma,y\co\tau]} M}{\Subst{z}{M}\rho}}
}
\\
\ib{\irule[\trule{\CR\Sigma}]
  {\WfTerm \Gamma %
    {\RecSigma {\rho[z\co\SigmaTy {x\co\sigma} \tau]}%
               {H[x\co\sigma,y\co\tau]} %
               {\SigmaPair {\tau[x\co\sigma]}{M}{N}}} %
    {\Subst z {\SigmaPair {\tau[x\co\sigma]} M N}\rho}};
  {\EqTerm \Gamma %
     {\RecSigma {\rho[z\co\SigmaTy {x\co\sigma}\tau]} {H[x\co\sigma,y\co\tau]} {\SigmaPair {\tau[x\co\sigma]} M N}}
     {\Substs{x\mapsto M, y\mapsto N}H}
     {\Subst z {\SigmaPair {\tau[x\co\sigma]} M N}\rho}}
}
\end{gather*}
\caption{Dependent Sum}
\end{figure}


\end{document}
%%% Local Variables:
%%% mode: LaTeX
%%% TeX-master: t
%%% End:
